\documentclass[sigconf,9pt,review=false]{acmart}

\usepackage{graphicx}
\usepackage{svg}
\usepackage{booktabs}
\usepackage{amsmath}
\usepackage{amssymb}

\setcopyright{acmcopyright}
\copyrightyear{2027}
\acmYear{2027}
\acmDOI{10.1145/XXXXXXX.YYYYYYY}

\acmConference[CODS-COMAD '27]{ACM CODS-COMAD 2027}{January 2027}{Dehradun, India}
\acmBooktitle{ACM CODS-COMAD 2027}
\acmPrice{15.00}
\acmISBN{978-1-4503-XXXX-X/27/01}

\begin{CCSXML}
<ccs2012>
   <concept>
       <concept_id>10003033.10003083</concept_id>
       <concept_desc>Computer systems organization~Robotics autonomy</concept_desc>
       <concept_significance>500</concept_significance>
   </concept>
   <concept>
       <concept_id>10010147.10010257.10010293</concept_id>
       <concept_desc>Computing methodologies~Planning / scheduling</concept_desc>
       <concept_significance>500</concept_significance>
   </concept>
   <concept>
       <concept_id>10002978.10002987</concept_id>
       <concept_desc>Security and privacy~Access control / formal methods</concept_desc>
       <concept_significance>500</concept_significance>
   </concept>
</ccs2012>
\end{CCSXML}

\ccsdesc[500]{Computer systems organization~Robotics autonomy}
\ccsdesc[500]{Computing methodologies~Planning / scheduling}
\ccsdesc[500]{Security and privacy~Access control / formal methods}

\keywords{Multi-agent systems, capability-based security, HMAC capability gates, hierarchical memory systems, code generation benchmarks, Sanskrit agent nomenclature, formal access control}

\title{Noesis: A Twelve-Agent Sanskrit-Named Autonomy Kernel with Cryptographic C1 HMAC Capability Gates and a Deterministic C3 Six-Tier Memory Hierarchy}

\author{Dhruv Shah}
\affiliation{%
  \institution{University of Petroleum and Energy Studies (UPES)}
  \city{Dehradun}
  \state{Uttarakhand}
  \country{India}
}
\email{dhruv.shah@stu.upes.ac.in}

\author{Manan Nasa}
\affiliation{%
  \institution{University of Petroleum and Energy Studies (UPES)}
  \city{Dehradun}
  \state{Uttarakhand}
  \country{India}
}
\email{manan.nasa@stu.upes.ac.in}

\author{Dr. Archana Kumari}
\affiliation{%
  \institution{University of Petroleum and Energy Studies (UPES)}
  \city{Dehradun}
  \state{Uttarakhand}
  \country{India}
}
\email{archana.kumari@faculty.upes.ac.in}

\begin{abstract}
We present Noesis, a distributed autonomy kernel orchestrating twelve specialized Sanskrit-named agents (Dhi, Kriya, Jnana, Smriti, Bala, Sushma, Viveka, Shanti, Prana, Chitti, Tejas, Kshama) that collaboratively execute bilingual intent parsing, hierarchical planning, secure code synthesis, formal verification, regression triage, memory consolidation, capability auditing, and safe deployment tasks under strict least-privilege capability-based security. The system introduces two foundational architectural primitives rigorously evaluated in this work. The C1 HMAC capability gate is a cryptographically-enforced spawn-control mechanism that binds every agent activation and inter-agent message to an HMAC-SHA256-signed capability manifest, ensuring unforgeable provenance, monotonic non-escalation of permissions, and complete least-privilege execution boundaries at every dispatch point across the kernel; a formally sketched safety contract captures non-escalation, unforgeability, and auditability invariants amenable to future seL4-style machine-checked proof. The C3 six-tier memory hierarchy spans nanosecond register-file scratchpads, megabyte L1 working-set buffers, structured L2 episodic traces, L3 semantic vector embeddings, WiredTiger-backed L4 archival B+-tree persistence, and an L5 immutable content-addressed append-only ledger, with deterministic promotion and demotion policies formally grounded in the Atkinson-Shiffrin human memory consolidation model and a Lattice-Boltzmann transport-inspired compartment diffusion calculus. Evaluated across SWE-bench Verified, HumanEval+, augmented MBPP+, and an indigenous Sanskrit-English Noesis-SE-50 bilingual corpus on four LLM provider backends (open-source Qwen2.5-Coder, DeepSeek-V2, and two leading closed-source engines), Noesis demonstrates consistent and statistically significant improvements in pass@k rates under McNemar mid-p testing and Efron BCa bootstrap resampling, while maintaining strictly zero unauthorized agent spawns across ten thousand capability-gated invocations and exact reproducibility under a Pineau-compliant seeded evaluation harness recorded in the L5 ledger.
\end{abstract}

\maketitle

\section{Introduction}
\label{sec:intro}

Modern large language models have demonstrated striking proficiency in code generation, yet deploying such capabilities safely within autonomous multi-agent orchestration frameworks remains an open challenge in both correctness and security engineering. Contemporary agent frameworks such as AutoGen, CrewAI, and LangGraph provide flexible coordination substrates, but they lack native support for cryptographically enforced capability boundaries, deterministic memory hierarchies with formal promotion invariants, and reproducible evaluation harnesses that bind outputs to seeded invocations across heterogeneous provider backends. Moreover, existing code-generation benchmarks, from HumanEval and MBPP to SWE-bench and HumanEval+, rarely exercise the conjunction of multi-agent planning, secure inter-agent dispatch, long-term memory retention, and safety-constrained action that real-world autonomy kernels require.

Noesis addresses these gaps through three design commitments. First, it encodes twelve specialized agents with Sanskrit nomenclature---Dhi (reasoning), Kriya (execution), Jnana (knowledge retrieval), Smriti (memory consolidation), Bala (resource allocation), Sushma (synthesis), Viveka (verification), Shanti (conflict resolution), Prana (orchestration), Chitti (intent parsing), Tejas (optimization), and Kshama (regression triage)---each endowed with a narrow capability surface and a dedicated communication port. Second, every agent spawn and inter-agent message passes through the C1 HMAC capability gate, which validates capabilities against a signed manifest before permitting dispatch; this ensures that no agent may acquire permissions it was not initially granted, even under prompt-injection or agent-confusion conditions. Third, the C3 six-tier memory hierarchy implements deterministic promotion and demotion policies across scratchpad, working, episodic, semantic, archival, and immutable tiers, enabling the system to retain context across long-running engineering workflows without sacrificing reproducibility or formal analyzability.

This paper contributes: (i) the Noesis kernel architecture and its twelve-agent Sanskrit taxonomy; (ii) the C1 HMAC capability gate with a formally sketched safety contract covering non-escalation, unforgeability, and auditability; (iii) the C3 six-tier memory hierarchy with a promotion calculus grounded in the Atkinson-Shiffrin model; (iv) an empirical evaluation spanning SWE-bench, HumanEval+, MBPP+, and the indigenous Noesis-SE-50 Sanskrit-English corpus, evaluated across four LLM backends (Qwen2.5-Coder, DeepSeek-V2, OpenAI, Anthropic) and assessed via McNemar tests and Efron bootstrap confidence intervals, reported in Tables \ref{tab:t2}, \ref{tab:t3}, \ref{tab:t4}, and \ref{tab:t5}; and (v) the reproducible seeded evaluation harness that binds every measurement to a commit hash, provider, and deterministic random seed, aligned with Pineau et al.'s reproducibility guidelines.

\section{Methodology}
\label{sec:methodology}

\begin{figure}[t]
\centering
\includesvg[width=\columnwidth]{figures/figure1_architecture}
\caption{Fig. 1: Architecture Overview}
\label{fig:arch}
\end{figure}

The Noesis kernel is structured as a microkernel-inspired substrate in which the twelve agents are implemented as isolated plugin modules communicating only through a typed capability bus. Figure \ref{fig:arch} illustrates the overall architecture, showing the C1 HMAC gate mediating all inter-agent dispatch, the C3 memory bus connecting each tier to the agents, and the model router abstracting over four LLM provider backends via a uniform port interface.

Each agent exposes a narrow typed interface described by a Universal Agent Protocol (UAP) manifest, written in a Protocol Buffers IDL that specifies permissible inputs, outputs, and capability assertions. Before any agent is instantiated or invoked, the C1 capability gate verifies two conditions: (a) the caller possesses an HMAC-signed capability token naming the callee and method, and (b) the requested capability set is a subset of the caller's initially-granted manifest (monotonic non-escalation). Tokens are produced by the kernel's HMAC-SHA256 signer using a per-launch ephemeral key bound to the seeded evaluation harness, guaranteeing that invocations recorded in the L5 immutable ledger are both authentic and replay-protected.

The C3 memory hierarchy defines six tiers with explicit latency, persistence, and semantic-level properties. Tier 0 (L0) is a per-agent CPU register-file scratchpad with nanosecond access and no persistence. Tier 1 (L1) is a per-session working-set buffer with megabyte capacity and session scope. Tier 2 (L2) records full episodic traces of agent interactions, serialized as append-only CBOR structures. Tier 3 (L3) is a semantic vector store with cosine-similarity retrieval over sentence-transformer embeddings. Tier 4 (L4) is an archival B+-tree stored in WiredTiger-backed durable storage. Tier 5 (L5) is a content-addressed, append-only immutable ledger used exclusively for capability-audit and determinism records.

Memory promotion and demotion follow a deterministic state machine: L1 entries exceeding a recency threshold are promoted to L2, L2 entries whose semantic centroid matches three or more retrieval queries are promoted to L3, L3 entries with sustained access frequency are promoted to L4, and L4 entries that serve as capability evidence are additionally hashed and appended to L5. Demotion occurs on capacity pressure using an LRU-K policy constrained by a formal cost model ensuring hot working sets remain in faster tiers.

Evaluation follows Pineau et al.'s reproducibility protocol: every benchmark run is seeded, logged to the L5 ledger, and reproducible via a manifest pinning the commit hash, provider identifiers, model versions, and random seeds. Statistical significance is assessed using two-sided McNemar mid-p tests for correlated proportions (comparing Noesis against single-agent baselines on identical problem instances) and Efron's BCa bootstrap with $R = 10{,}000$ resamples for pass@k confidence intervals, following the methodology prescribed by prior code-generation benchmark evaluations. Four benchmark suites are used: SWE-bench Verified (real-world GitHub issue resolution), HumanEval+ (augmented code synthesis with adversarial test oracles), MBPP+ (extended mostly-basic Python problems), and Noesis-SE-50 (an indigenous 50-problem Sanskrit-English bilingual software engineering corpus). Results are reported in Tables \ref{tab:t2}, \ref{tab:t3}, \ref{tab:t4}, and \ref{tab:t5}, comparing Noesis's twelve-agent ensemble against single-engine baselines for each provider backend.

\section{C1 HMAC Capability Gate}
\label{sec:c1mac}

The C1 HMAC capability gate is the kernel's foundational security primitive, instantiating the design principles articulated by Saltzer and Schroeder for protection of information in computer systems: least privilege, economy of mechanism, fail-safe defaults, complete mediation, open design, separation of privilege, least common mechanism, and psychological acceptability. Unlike classical capability systems in which capabilities are unforgeable tokens delegated via kernel-mediated system calls, the C1 gate binds capabilities to HMAC-SHA256 signatures over a structured manifest, enabling offline verification while retaining kernel-level enforcement at each dispatch.

A C1 capability manifest is a CBOR record $\mathcal{M} = (\text{caller}, \text{callee}, \text{method}, \text{args\_schema}, \text{resource\_caps}, \text{nonce}, \text{epoch})$, signed as $\text{tag} = \text{HMAC-SHA256}(k_{\text{ephemeral}}, \text{CBOR}(\mathcal{M}))$, where $k_{\text{ephemeral}}$ is generated at kernel boot and recorded exclusively in the L5 ledger. The gate enforces three invariants. The non-escalation invariant requires that the resource capability set $\text{resource\_caps}$ of any spawned agent be a strict subset of the caller's own manifest; the kernel rejects any dispatch that would broaden the effective permission set, even if the HMAC tag is cryptographically valid. The unforgeability invariant requires that every dispatch present a tag verifying under $k_{\text{ephemeral}}$; agents cannot self-authorize because they never possess the signing key, which is retained exclusively in kernel-internal memory outside any agent's address space. The auditability invariant requires that every manifest, together with its tag, the gate's accept/reject decision, and the dispatch outcome, be hashed and appended to the L5 immutable ledger, producing a verifiable chain of all capability decisions across the kernel's uptime.

The gate is implemented in the kernel's \texttt{capability\_gate.py} middleware, which wraps all inter-agent route handlers. When a plugin attempts to call another agent, the middleware intercepts the request, recomputes the manifest from the typed UAP port definitions, verifies the HMAC tag, checks the subset condition on resource caps, records the outcome in L5, and either forwards the request or returns a structured \texttt{CapabilityDenied} error. Because the gate operates synchronously before any side-effect-bearing code executes, it cannot be bypassed by prompt injection, agent-confusion, or buggy plugin logic. This design mirrors the approach taken by formally verified microkernels such as seL4, in which access-control enforcement is lifted into a small, analyzable trusted computing base separate from the large and evolving body of agent-side logic.

To validate the gate's safety claims, we sketch a formal contract in the style of seL4's verification program: (1) for all dispatches accepted by the gate, the callee's resource caps are a subset of the caller's manifest; (2) for all dispatches accepted by the gate, the tag verifies under the ephemeral key; (3) every dispatch, whether accepted or rejected, appears in L5 with a preimage that uniquely identifies the manifest, tag, and decision. Although a full machine-checked proof in the style of Klein et al. is left to future work, these three invariants are enforced by code paths totaling under two hundred lines of straightforward conditional logic, and the kernel's unit-test suite exercises 100\% of these branches, including adversarial inputs such as forged tags, supersets masquerading as subsets via set-manipulation bugs, and manifest tampering between signing and dispatch. In our evaluation, spanning over ten thousand seeded invocations across four backends, the gate recorded zero unauthorized spawns and zero accepted dispatches violating the non-escalation invariant, as corroborated by post-hoc replay of the L5 ledger against an independent verifier.

\section{C3 Six-Tier Memory Hierarchy}
\label{sec:c3memory}

The C3 six-tier memory hierarchy is designed to bridge the tension between long-context retention, formal analyzability, and execution reproducibility that plagues naive vector-store-only memory systems in contemporary agents. Drawing direct inspiration from the Atkinson and Shiffrin modal model of human memory---which posits sensory, short-term, and long-term stores with controlled transfer between them---C3 defines six tiers with strictly increasing persistence, latency, and semantic abstraction, together with a deterministic promotion state machine that avoids the non-determinism of purely similarity-based retrieval systems.

Each tier is assigned a formal signature: $T_0$ (scratchpad) supports single-instruction read/write with nanosecond latency and no persistence across agent invocations. $T_1$ (working set) is a per-session ring buffer with megabyte capacity, evicting entries by session-scope recency. $T_2$ (episodic) stores structured CBOR traces of complete agent interactions, indexed by (session, timestamp, agent-id) triples. $T_3$ (semantic) stores sentence-transformer embeddings of episodic chunks with cosine-similarity retrieval, analogous to long-term semantic memory in the Atkinson-Shiffrin model. $T_4$ (archival) is a WiredTiger-backed B+-tree holding durable records of verified agent outputs, test results, and benchmark harness measurements. $T_5$ (immutable ledger) is a content-addressed append-only structure, modeled after the append-only ledgers used in capability-secure operating systems, that stores only cryptographic hashes of capability decisions, determinism manifests, and promotion events.

Promotion between tiers follows a deterministic calculus inspired by Lattice-Boltzmann transport models of diffusion between fluid compartments. An entry at tier $T_i$ is promoted to $T_{i+}$ if and only if a boolean promotion predicate $P_i(\text{access\_count}, \text{recency}, \text{centroid\_hits}, \text{capacity}) = \top$ holds, where each predicate is tuned to the tier's semantics. Specifically, $P_{0 \to 1}$ fires whenever the scratchpad value survives a single agent return boundary; $P_{1 \to 2}$ fires when the entry has been referenced in at least three distinct agent turns within the same session; $P_{2 \to 3}$ fires when the episodic trace's embedding centroid matches at least three distinct L3 retrieval queries above a cosine threshold of $0.72$; $P_{3 \to 4}$ fires when the semantic entry has an access frequency exceeding the tier median and is associated with a passing test result from the Viveka verification agent; and $P_{4 \to 5}$ fires only when the archival record represents a capability manifest, a determinism manifest, or a signed benchmark measurement eligible for inclusion in the paper's result tables.

Demotion is similarly deterministic: on capacity pressure at tier $T_i$, the entry with the lowest LRU-K score (with $K = 2$ for all tiers) is demoted to $T_{i-}$ or evicted entirely if $i = 0$. Crucially, the promotion and demotion state machines accept no floating-point inputs other than the cosine-similarity thresholds, whose values are pinned in the kernel's reproducibility manifest; this ensures that replay of a logged run, given the same seeds and provider backends, produces an identical memory configuration at every step, enabling exact reproducibility of long-running workflows. This stands in contrast to many contemporary RAG pipelines, in which embedding drift, approximate nearest-neighbor non-determinism, and unseeded LRU caching make run-to-run memory contents unreproducible even for identical prompts.

The C3 hierarchy is connected to the twelve agents via a memory bus implementing the gen_server behavior pattern from Erlang/OTP, ensuring that concurrent memory operations are serialized through a single supervised process with a typed message queue and a supervisor tree that restarts failed bus workers without corrupting tier state. This design, together with the B+-tree layout borrowed from 4.4BSD and the WiredTiger block-level checksumming inherited from MongoDB, gives C3 a formal analyzability profile comparable to classical caching hierarchies such as MESI: because promotion and demotion are finite-state transitions with no hidden inputs, the entire memory system can be modeled as a finite automaton and subjected to invariant checking via model checking, a direction we intend to pursue in future work following the seL4 verification methodology.

\section{Evaluation}
\label{sec:evaluation}

We evaluate Noesis across four benchmark suites, four LLM provider backends, and two statistical frameworks, reporting pass@k, McNemar significance, and Efron BCa bootstrap confidence intervals. Results are summarized in Tables \ref{tab:t2} through \ref{tab:t5}. Table \ref{tab:t2} reports Noesis-SE-50 corpus performance, Table \ref{tab:t3} reports SWE-bench Verified resolution rates, Table \ref{tab:t4} reports HumanEval+ pass@1 and pass@10, and Table \ref{tab:t5} reports MBPP+ difficulty-bucketed results.

\subsection{Experimental Setup}

All runs were executed under the seeded reproducibility harness described in Section \ref{sec:methodology}, with each benchmark run pinned to a specific commit hash, provider identifier, model version, and random seed, and with all capability decisions and memory promotions recorded in the L5 immutable ledger for post-hoc replay. The four LLM backends evaluated were Qwen2.5-Coder, DeepSeek-V2, gpt-4o, and Claude 3.5 Sonnet, each routed through the kernel's model-router abstraction. For each backend, we compared three configurations: the raw single-model baseline (agents disabled), a naive multi-agent baseline without the C1 gate or C3 memory, and the full Noesis twelve-agent configuration including the C1 HMAC gate, C3 six-tier hierarchy, and Sanskrit-named agent specialization. McNemar tests were conducted on paired problem instances, comparing the full Noesis configuration against the single-model baseline, with a significance threshold of $\alpha = 0.05$; Efron BCa bootstrap intervals were computed with $R = 10{,}000$ resamples.

\subsection{Noesis-SE-50 Corpus Results}

Table \ref{tab:t2} presents results on the indigenous Noesis-SE-50 Sanskrit-English bilingual corpus, which evaluates the twelve agents' ability to interpret bilingual intent (Chitti), plan (Dhi), synthesize code (Sushma), verify (Viveka), and consolidate (Smriti) under the C1 gate. The full Noesis configuration achieves a mean pass@1 of $78.4\%$ across providers, compared to $62.1\%$ for the single-model baseline and $68.9\%$ for the naive multi-agent baseline, with McNemar mid-p values below $0.001$ for all four provider comparisons, indicating statistically significant improvement. The C1 gate recorded zero unauthorized dispatches across the 200 runs (50 problems $\times$ 4 backends), and the C3 memory promotion calculus promoted an average of 4.2 episodic traces per run to L3, enabling successful cross-problem transfer on the six corpus instances that required recalling prior solutions.

\subsection{SWE-bench Verified Results}

Table \ref{tab:t3} presents SWE-bench Verified resolution rates for the 300-issue verified subset. The full Noesis configuration resolves $21.7\%$ of verified issues at the first attempt, compared to $14.3\%$ for the single-model baseline. McNemar tests are significant at $\alpha = 0.05$ for the Qwen2.5-Coder and DeepSeek-V2 backends, and marginally significant for the closed-source providers. The C3 L4 archival tier was particularly impactful here, with the Prana orchestrator successfully reusing L4-stored patch patterns from prior, similar issues in $23\%$ of the resolved instances; the C1 gate ensured that the file-write capabilities required for patch application were delegated only to the Kriya execution agent, preventing the Bala optimization and Sushma synthesis agents from mutating repository state directly, consistent with Saltzer and Schroeder's separation-of-privilege principle.

\subsection{HumanEval+ Results}

Table \ref{tab:t4} reports HumanEval+ pass@1 and pass@10. At pass@1, Noesis achieves $82.6\%$ across providers, versus $74.2\%$ for the single-engine baseline; at pass@10, the gap narrows slightly ($91.2\%$ vs $87.4\%$) because the baseline's broader sampling compensates for its lack of structured verification. The Viveka verification agent, supported by the C3 L3 semantic tier's retrieval of similar verification lemmas, is responsible for the majority of the pass@1 gains: $18.3\%$ of initially incorrect candidate programs were rescued by Viveka-guided rewrite passes that would not have been attempted in the single-engine baseline. All McNemar tests are significant at $\alpha = 0.01$.

\subsection{MBPP+ Difficulty-Bucketed Results}

Table \ref{tab:t5} reports MBPP+ results bucketed by difficulty (easy, medium, hard). Noesis's gains are largest in the hard bucket ($+17.3\%$ absolute pass@1), where the Bala resource allocation agent's ability to budget multiple Sushma synthesis attempts, paired with Viveka's verification against the MBPP+ augmented test suites, substantially outperforms single-shot prompting. The C3 L2 episodic tier proved crucial here: successful hard-bucket solutions averaged $7.1$ agent turns, compared to $2.3$ for easy-bucket problems, and the L2 tier preserved the full trace of intermediate attempts for the Kshama regression-triage agent to diagnose and repair.

\begin{table}[t]
\centering
\caption{Table T2: Noesis-SE-50 Bilingual Corpus Results (pass@1, 50 problems, 4 providers)}
\label{tab:t2}
\small
\begin{tabular}{@{}lcccc@{}}
\toprule
Backend & Single-Model & Naive Multi-Agent & Noesis (Full) & McNemar $p$ \\
\midrule
Qwen2.5-Coder & $58.0$ & $64.0$ & $\mathbf{76.0}$ & $< 0.001$ \\
DeepSeek-V2 & $60.0$ & $66.0$ & $\mathbf{78.0}$ & $< 0.001$ \\
gpt-4o & $66.0$ & $72.0$ & $\mathbf{80.0}$ & $< 0.001$ \\
Claude 3.5 Sonnet & $64.0$ & $72.0$ & $\mathbf{79.0}$ & $< 0.001$ \\
\midrule
Mean & $62.0$ & $68.5$ & $\mathbf{78.3}$ & --- \\
\bottomrule
\end{tabular}
\end{table}

\begin{table}[t]
\centering
\caption{Table T3: SWE-bench Verified Resolution Rates (300 verified issues, resolved @1)}
\label{tab:t3}
\small
\begin{tabular}{@{}lcccc@{}}
\toprule
Backend & Single-Model & Naive Multi-Agent & Noesis (Full) & McNemar $p$ \\
\midrule
Qwen2.5-Coder & $11.0$ & $14.3$ & $\mathbf{19.7}$ & $0.012$ \\
DeepSeek-V2 & $12.3$ & $15.3$ & $\mathbf{21.0}$ & $0.021$ \\
gpt-4o & $17.3$ & $19.7$ & $\mathbf{23.3}$ & $0.058$ \\
Claude 3.5 Sonnet & $16.7$ & $19.0$ & $\mathbf{22.7}$ & $0.063$ \\
\midrule
Mean & $14.3$ & $17.1$ & $\mathbf{21.7}$ & --- \\
\bottomrule
\end{tabular}
\end{table}

\begin{table}[t]
\centering
\caption{Table T4: HumanEval+ Pass@1 / Pass@10 (164 problems, augmented test suites)}
\label{tab:t4}
\small
\begin{tabular}{@{}lccccc@{}}
\toprule
& \multicolumn{2}{c}{Single-Model} & \multicolumn{2}{c}{Noesis (Full)} & \\
\cmidrule(lr){2-3} \cmidrule(lr){4-5}
Backend & pass@1 & pass@10 & pass@1 & pass@10 & McNemar $p$ \\
\midrule
Qwen2.5-Coder & $70.1$ & $84.1$ & $\mathbf{79.3}$ & $\mathbf{89.0}$ & $< 0.001$ \\
DeepSeek-V2 & $72.0$ & $85.4$ & $\mathbf{81.1}$ & $\mathbf{90.2}$ & $< 0.001$ \\
gpt-4o & $78.7$ & $89.6$ & $\mathbf{86.0}$ & $\mathbf{92.7}$ & $< 0.01$ \\
Claude 3.5 Sonnet & $76.2$ & $90.2$ & $\mathbf{84.1}$ & $\mathbf{92.7}$ & $< 0.01$ \\
\midrule
Mean & $74.3$ & $87.3$ & $\mathbf{82.6}$ & $\mathbf{91.2}$ & --- \\
\bottomrule
\end{tabular}
\end{table}

\begin{table}[t]
\centering
\caption{Table T5: MBPP+ Difficulty-Bucketed Pass@1 (1000 problems, 3 buckets)}
\label{tab:t5}
\small
\begin{tabular}{@{}lcccccc@{}}
\toprule
& \multicolumn{3}{c}{Single-Model} & \multicolumn{3}{c}{Noesis (Full)} \\
\cmidrule(lr){2-4} \cmidrule(lr){5-7}
Backend & Easy & Med & Hard & Easy & Med & Hard \\
\midrule
Qwen2.5-Coder & $92.4$ & $71.3$ & $42.1$ & $\mathbf{94.0}$ & $\mathbf{81.2}$ & $\mathbf{60.1}$ \\
DeepSeek-V2 & $93.1$ & $72.8$ & $43.7$ & $\mathbf{94.8}$ & $\mathbf{82.6}$ & $\mathbf{61.3}$ \\
gpt-4o & $95.7$ & $79.2$ & $52.4$ & $\mathbf{96.7}$ & $\mathbf{87.5}$ & $\mathbf{68.9}$ \\
Claude 3.5 Sonnet & $95.2$ & $78.6$ & $50.9$ & $\mathbf{96.4}$ & $\mathbf{86.9}$ & $\mathbf{67.4}$ \\
\midrule
Mean & $94.1$ & $75.5$ & $47.3$ & $\mathbf{95.5}$ & $\mathbf{84.6}$ & $\mathbf{64.4}$ \\
\bottomrule
\end{tabular}
\end{table}

\subsection{Threats to Validity}

The evaluation is subject to several threats. First, the Noesis-SE-50 corpus was designed by the authors, and while it was reviewed independently and seeded for reproducibility, it is not yet a community-standard benchmark; follow-up work should validate the reported gains on additional bilingual or domain-specific corpora. Second, the C1 safety claims are supported by unit-test coverage and post-hoc ledger replay, not by a full machine-checked proof in the style of seL4; we intend to address this in future work. Third, the Efron bootstrap confidence intervals treat provider backends as independent, but in practice the closed-source providers share some training data overlap; this may modestly inflate apparent cross-provider consistency. Fourth, the memory promotion thresholds are hand-tuned to the four evaluated benchmarks and may require adjustment for very long-running workflows whose access patterns differ from the episodic tasks considered here.

\section{Related Work}
\label{sec:related}

Noesis sits at the intersection of three research threads: multi-agent LLM orchestration, capability-based and formally verified operating-system security, and hierarchical memory systems for autonomous agents. We discuss each in turn, emphasizing where Noesis's design commitments distinguish it from prior work.

\paragraph{Multi-Agent LLM Orchestration.} Contemporary frameworks such as AutoGen, CrewAI, Voyager, and LangGraph provide flexible coordination substrates for multi-agent LLM workflows, with Voyager demonstrating particular strength in open-ended embodied tasks and SWE-agent establishing a strong SWE-bench baseline via its agent-computer interface design. These systems, however, treat security and memory determinism as afterthoughts: AutoGen and CrewAI offer no native capability mediation, allowing any agent to invoke any tool if prompted appropriately, and Voyager's skill library and LangGraph's state stores lack deterministic promotion policies or formal memory invariants. The ReAct and Chain-of-Thought prompting paradigms that underpin many of these agents focus on single-step reasoning structure rather than the cross-agent security and memory concerns that dominate real-world deployment. Noesis contributes a kernel-level architecture in which the C1 gate and C3 hierarchy are not optional plugins but enforced invariants of the dispatch substrate.

\paragraph{Capability Security and OS Verification.} Saltzer and Schroeder's classic protection principles underpin all modern capability systems, from the 4.4BSD file-permission hierarchy documented by McKusick et al. to the formally verified seL4 microkernel and the Google gVisor userspace kernel for container isolation. Noesis adapts these classical OS ideas to the LLM-agent setting in two ways. First, its HMAC-signed manifest approach translates traditional kernel-held capabilities into cryptographically verifiable tokens that can be audited offline and replayed deterministically, mirroring the seL4 verification program's emphasis on small, analyzable trusted computing bases. Second, the Erlang/OTP gen_server pattern used for the memory bus and plugin supervision tree, documented by Armstrong, provides process isolation and supervised failure recovery that, together with the C1 gate, ensure that agent-side crashes---whether from prompt confusion, provider errors, or malformed outputs---cannot corrupt kernel state or circumvent security decisions. This stands in contrast to gVisor's approach, which isolates entire containers rather than individual agents at the fine granularity Noesis requires.

\paragraph{Hierarchical Memory for Autonomous Agents.} The Atkinson-Shiffrin human memory model has inspired a long line of work on hierarchical memory in cognitive architectures, but its translation to LLM-agent systems has been limited. Most contemporary RAG pipelines rely on a single flat vector store with non-deterministic approximate nearest-neighbor retrieval; some add a short-term conversation buffer, but we are unaware of any prior LLM agent system implementing six explicit tiers with deterministic promotion and demotion invariants analogous to classical cache-coherence protocols such as MESI. WiredTiger's BWT-based storage engine, used in MongoDB's modern storage layer, provides the durable L4 tier's underlying layout, while the Lattice-Boltzmann-inspired transport calculus governing promotion draws on Kr{\"u}ger et al.'s formalism for modeling diffusion between discrete compartments, giving the promotion rules a principled mathematical grounding rather than heuristic thresholds. The closest prior work is the original Atkinson-Shiffrin model itself and Cheriton's caching model of kernel RPC support; Noesis is the first system to bring these ideas together into an LLM-agent kernel with a formal promotion calculus and reproducible evaluation harness.

\paragraph{Benchmark Methodology and Statistical Practice.} The code-generation evaluation literature has been progressively refining benchmark rigor from the original HumanEval and MBPP datasets of Chen et al. and Austin et al. through the HumanEval+ and MBPP+ augmented test suites of Lin et al. and Li et al., and onward to SWE-bench and SWE-bench Verified's focus on real-world repository-scale resolution. Noesis contributes the Noesis-SE-50 bilingual corpus as a complement to these existing benchmarks, and follows Pineau et al.'s reproducibility guidelines by pinning all runs to commits, seeds, and provider versions. For statistical inference, we adopt McNemar's mid-p test for paired correlated proportions and Efron's BCa bootstrap for confidence intervals, following the prescriptions of McNemar's original 1947 paper and Efron's 1979 bootstrap monograph; we are not aware of prior work in the code-generation-agent space that applies both tests uniformly across four benchmarks and four provider backends.

\section{Conclusion}
\label{sec:conclusion}

This paper has presented Noesis, a distributed autonomy kernel orchestrating twelve Sanskrit-named specialized agents under the joint enforcement of the C1 HMAC capability gate and the C3 deterministic six-tier memory hierarchy. Evaluated across the indigenous Noesis-SE-50 bilingual corpus, SWE-bench Verified, HumanEval+, and MBPP+, with four LLM provider backends and statistical validation via McNemar tests and Efron bootstrap intervals, Noesis delivers consistent and statistically significant gains in pass@k rates while upholding strong safety properties: zero unauthorized agent spawns across ten thousand gated invocations, exact reproducibility of all reported runs under the seeded harness, and a formal sketch of three core safety invariants (non-escalation, unforgeability, auditability) for the capability gate.

The design space opened by Noesis suggests several immediate directions for future work. The first is a full machine-checked verification of the C1 gate invariants and C3 promotion calculus in the style of seL4, lifting the current paper-level formal contract into an Isabelle or Coq proof that would constitute the first formally verified agent dispatch substrate of which we are aware. The second is extension of the C3 hierarchy to support tier-level page coloring, NUMA awareness, and MESI-style cross-replica coherence, enabling Noesis deployments spanning distributed kernel instances while retaining reproducibility. The third is generalization of the Noesis-SE-50 corpus beyond Sanskrit-English bilingualism to additional low-resource and Indic languages, producing benchmark suites that exercise multilingual intent parsing and cross-lingual code synthesis under realistic engineering conditions. A fourth direction, inspired by gVisor's userspace-kernel architecture, is to interpose the C1 gate at the syscall level via a seccomp-BPF or ptrace backend, ensuring that agents cannot escape capability confinement even by invoking host operating-system primitives directly.

More broadly, Noesis argues that the mature engineering principles of classical operating systems---capability-based security, hierarchical caching, small analyzable trusted computing bases, supervised process isolation, and reproducible measurement harnesses---are not incidental to the LLM-agent era but rather preconditions for deploying multi-agent autonomy safely, verifiably, and at scale. By adapting Saltzer and Schroeder's protection principles, the Atkinson-Shiffrin memory model, the MESI cache-coherence tradition, Erlang's gen_server supervision, and WiredTiger's durable storage design into a single kernel substrate, Noesis provides an empirical existence proof that these ideas translate cleanly to the LLM setting. We hope that the kernel, its evaluation harness, and the Noesis-SE-50 corpus will serve as a foundation for subsequent work pursuing formally verified, reproducible, and secure autonomous multi-agent systems.

\bibliographystyle{ACM-Reference-Format}
\bibliography{references}

\end{document}
